\par\Needspace{6\baselineskip}\medskip
\hypertarget{native-208-article154}{}\pdfbookmark[2]{154.}{article-154}
\begin{center}154.\end{center}
Ne quidquam desiderandum relinquatur, adhuc \textit{tertiam hypothesin} extruemus, ubi rursus valores ipsarum $P',Q'$ in hypothesi secunda erutos tamquam valores ipsarum $P,Q$ adoptabimus. Statuendo itaque
\[
\begin{aligned}
x&=\log P=0{,}079\,0167\\
y&=\log Q=8{,}547\,6110,
\end{aligned}
\]
praecipua calculi momenta haec inveniuntur:
\begin{center}\small\setlength{\tabcolsep}{6pt}
\begin{tabular}{@{}r|r@{}}
$\omega=13^{0}\,15'\,38''{,}39$ & $\zeta''=210^{0}\,8'\,25''{,}65$\\[.5em]
$\omega+\sigma=13\quad38\quad51{,}51$ & $\log r=0{,}330\,7640$\\[.5em]
$\log Qc\sin\omega=0{,}598\,9542$ & $\log r''=0{,}322\,2239$\\[.5em]
$z=14^{0}\,33'\,19''{,}50$ & $\tfrac12(u''+u)=205^{0}\,22'\,14''{,}57$\\[.5em]
$\log r'=0{,}325\,9878$ & $\tfrac12(u''-u)=-3\quad14\quad4{,}78$\\[.5em]
$\log\dfrac{n'r'}{n}=0{,}667\,5154$ & $2f'=7\quad34\quad53{,}73$\\[.7em]
$\log\dfrac{n'r'}{n''}=0{,}588\,4987$ & $2f=3\quad29\quad0{,}39$\\[.7em]
$\zeta=203^{0}\,16'\,38''{,}41$ & $2f''=4\quad5\quad53{,}34$
\end{tabular}
\end{center}

Omnes hi numeri ab iis quos hypothesis secunda suppeditaverat tam parum differunt, ut certo concludere liceat, hypothesin tertiam nulla amplius correctione indigere\footnote{Si calculus perinde ut in hypothesibus antecedentibus ad finem perduceretur, prodiret $X=0$, $Y=-0{,}000\,0003$, qui valor tamquam evanescens considerandus est, et vix supra incertitudinem figurae decimali ultimae semper inhaerentem exsurgit.}. Progredi itaque licet ad ipsam elementorum determinationem ex $2f',r,r'',\theta'$, quam huc transscribere supersedemus, quum iam supra art.~97 exempli loco in extenso allata sit. Nihil itaque superest, nisi ut positionem plani orbitae ad normam art.~149 computemus, epochamque ad initium anni 1805 transferamus. Calculus ille superstruendus est numeris sequentibus:
\SourceLeaf{209}
\[
\begin{aligned}
AD'-\zeta&=9^{0}\,55'\,51''{,}41\\
\tfrac12(\gamma+u)&=202\quad18\quad13{,}855\\
\tfrac12(\gamma-u)&=-6\quad18\quad5{,}495,
\end{aligned}
\]
unde derivamus
\[
\begin{aligned}
\tfrac12(g+h)&=196^{0}\,43'\,14''{,}62\\
\tfrac12(g-h)&=-4\quad37\quad24{,}41\\
\tfrac12i&=6\quad33\quad22{,}05.
\end{aligned}
\]
Fit igitur $h=201^{0}\,20'\,39''{,}03$, adeoque $\AscendingNode=l-h=171^{0}\,7'\,48''{,}73$; porro $g=192^{0}\,5'\,50''{,}21$, et proin, quum anomalia vera pro observatione prima in art.~97 inventa sit $=310^{0}\,55'\,29''{,}64$, distantia perihelii a nodo ascendente in orbita $=241^{0}\,10'\,20''{,}57$, longitudoque perihelii $=52^{0}\,18'\,9''{,}30$; denique inclinatio orbitae $=13^{0}\,6'\,44''{,}10$. --- Si ad eundem calculum a loco tertio proficisci malumus, habemus
\[
\begin{aligned}
A''D'-\zeta''&=24^{0}\,18'\,35''{,}25\\
\tfrac12(\gamma''+u'')&=196\quad24\quad54{,}98\\
\tfrac12(\gamma''-u'')&=-5\quad43\quad14{,}81.
\end{aligned}
\]
Hinc elicitur
\[
\begin{aligned}
\tfrac12(g''+h'')&=211^{0}\,24'\,32''{,}44\\
\tfrac12(g''-h'')&=-11\quad43\quad48{,}49\\
\tfrac12i&=6\quad33\quad22{,}05,
\end{aligned}
\]
atque hinc longitudo nodi ascendentis $=l''-h''=171^{0}\,7'\,48''{,}72$, longitudo perihelii $=52^{0}\,18'\,9''{,}30$, inclinatio orbitae $=13^{0}\,6'\,44''{,}10$, prorsus eaedem ut ante.

Intervallum temporis ab observatione ultima usque ad initium anni 1805 est dierum $64{,}614\,102$; cui respondet motus heliocentricus medius $53293''{,}65=14^{0}\,48'\,13''{,}65$; hinc fit epocha anomaliae mediae pro initio anni 1805 in meridiano Parisino $=349^{0}\,34'\,12''{,}38$, atque epocha longitudinis mediae $=41^{0}\,52'\,21''{,}68$.
